\subsection{Comparison of two cohomotopers}

Consider a diagram

\[
\begin{array}{c} \mathrm{H} \xrightarrow {\varphi} \mathrm{G} \\ u \Big \downarrow \qquad \qquad \qquad \qquad \qquad \qquad \qquad \qquad \qquad \Big \downarrow v \\ \mathrm{H} ^ {\prime} \xrightarrow {\varphi^ {\prime}} \mathrm{G} ^ {\prime} \\ \psi^ {\prime} \end{array}\tag{4}
\]

where H, H', G, G' are groupoids and u, v, \(\varphi\) , \(\psi\) , \(\varphi'\) , \(\psi'\) are morphisms of groupoids such that \(v \circ \varphi = \varphi' \circ u\) and \(v \circ \psi = \psi' \circ u\) .

Denote by \(\alpha\) the canonical morphism from G to the cohomotoper \(\mathrm{Coh}(\varphi, \psi)\) and \(h\) the canonical homotopy connecting \(\alpha \circ \varphi\) to \(\alpha \circ \psi\) ; define \(\alpha'\) and \(h'\) analogously. Then, \(\alpha' \circ v\) is a groupoid morphism from G to \(\mathrm{Coh}(\varphi', \psi')\) and \(h' \circ \mathrm{Som}(u)\) is a homotopy connecting \(\alpha' \circ \varphi' \circ u\) to \(\alpha' \circ \psi' \circ u\) , that is, \(\alpha' \circ v \circ \varphi\) to \(\alpha' \circ v \circ \psi\) . By the universal property of cohomotopers (II, p. 185, prop. 3), there exists a unique morphism of groupoids \(w\) from \(\mathrm{Coh}(\varphi, \psi)\) to \(\mathrm{Coh}(\varphi', \psi')\) such that \(w \circ \alpha = \alpha' \circ v\) and \(\mathrm{Fl}(w) \circ h = h' \circ \mathrm{Som}(u)\) . We have in particular extended diagram (4) to a diagram

\[
\begin{array}{c} \mathrm{H} \xrightarrow [ \psi ]{\varphi} \mathrm{G} \xrightarrow {\alpha} \mathrm{Coh} (\varphi , \psi) \\ u \Biggl \downarrow \qquad \qquad \qquad \Biggl \downarrow v \qquad \qquad \Biggl \downarrow w \\ \mathrm{H} ^ {\prime} \xrightarrow [ \psi^ {\prime} ]{\varphi^ {\prime}} \mathrm{G} ^ {\prime} \xrightarrow {\alpha^ {\prime}} \mathrm{Coh} (\varphi^ {\prime}, \psi^ {\prime}) \end{array}\tag{5}
\]

in which the second square is commutative.

THEOREM 1. --- Assume the following hypotheses:

\begin{enumerate}
\def\labelenumi{(\roman{enumi})}
\item
  the morphism of groupoids v is a homotopism;
\item
  the map \(\operatorname{Orb}(u)\colon\operatorname{Orb}(\mathrm{H})\to\operatorname{Orb}(\mathrm{H}')\) , deduced from u by passing to orbits, is bijective;
\item
  there exists in each orbit of H a point a such that the homomorphism \(u_a \colon \mathrm{H}_a \to \mathrm{H}_{u(a)}'\) is surjective.
\end{enumerate}

Then, the morphism of groupoids \(w\colon \operatorname{Coh}(\varphi,\psi)\to\operatorname{Coh}(\varphi',\psi')\) is a homotopism.

Let G'' be the groupoid deduced from G' by contracting the arrows of a maximal oriented forest. The canonical morphism \(v'\colon G' \to G''\) is a homotopism (II, p. 184, corollary 2 of prop. 1). The two diagrams

\[
\begin{array}{c c c} \mathrm{H} \xrightarrow [ \psi ]{\varphi} \mathrm{G} & & \mathrm{H} ^ {\prime} \xrightarrow [ \psi^ {\prime} ]{\varphi^ {\prime}} \mathrm{G} ^ {\prime} \\ u \Big \downarrow & \Big \downarrow_ {v ^ {\prime} \circ v} & \text {and} \quad \mathrm{Id} _ {\mathrm{H} ^ {\prime}} \Big \downarrow \\ \mathrm{H} ^ {\prime} \xrightarrow [ v ^ {\prime} \circ \psi^ {\prime} ]{v ^ {\prime} \circ \varphi^ {\prime}} \mathrm{G} ^ {\prime} & & \mathrm{H} ^ {\prime} \xrightarrow [ v ^ {\prime} \circ \psi^ {\prime} ]{v ^ {\prime} \circ \varphi^ {\prime}} \mathrm{G} ^ {\prime \prime} \end{array}
\]

give rise to morphisms of groupoids \(w_1' \colon \mathrm{Coh}(\varphi, \psi) \to \mathrm{Coh}(v' \circ \varphi', v' \circ \psi')\) and \(w_2' \colon \mathrm{Coh}(\varphi', \psi') \to \mathrm{Coh}(v' \circ \varphi', v' \circ \psi')\) ; we have \(w_1' = w_2' \circ w\) . It therefore suffices to prove that \(w_1'\) and \(w_2'\) are homotopisms. Since the maps \(\mathrm{Som}(v' \circ v) \colon \mathrm{Som}(\mathrm{G}) \to \mathrm{Som}(\mathrm{G}'' )\) and \(\mathrm{Som}(v') \colon \mathrm{Som}(\mathrm{G}') \to \mathrm{Som}(\mathrm{G}'' )\) are surjective, it consequently suffices to prove the theorem under the additional hypothesis that the map \(\mathrm{Som}(v)\) is surjective, a hypothesis that we shall make throughout the rest of the proof.

Let us prove successively the following assertions:

- The map Orb(w) is bijective;

- For every vertex \(a\) of G, the homomorphism \(w_{a}\) is surjective;

- For every vertex \(a\) of G, the homomorphism \(w_{a}\) is injective.

\begin{enumerate}
\def\labelenumi{\alph{enumi})}
\item
  By hypothesis, the map \(\operatorname{Orb}(u)\) is bijective; the same holds for the map \(\operatorname{Orb}(v)\) according to II, p. 182, prop. 1, since v is a homotopism. The quiver morphism from the frame of the pair \((\varphi,\psi)\) to that of the pair \((\varphi',\psi')\) defined by the maps \(\operatorname{Orb}(u)\) and \(\operatorname{Orb}(v)\) is therefore an isomorphism. In particular, the map deduced from it by passing to connected components is bijective. Prop. 4 of II, p. 185 then implies that the map \(\operatorname{Orb}(w)\) is bijective.
\item
Let \(f'\) be an arrow of \(G'\) , denote by \(a'\) its origin and \(b'\) its target. Since the map \(\operatorname{Som}(v)\) is surjective, there exist vertices \(a\) and \(b\) in G such that \(a' = v(a)\) and \(b' = v(b)\) . Since the morphism v is a homotopism, there exists an arrow \(f\) of G connecting \(a\) to \(b\) , as well as an element \(g \in G_{a}\) such that \(v(g) = f'v(f)^{-1}\) (II, p. 182, prop. 1), whence \(f' = v(gf)\) . This proves that the map \(\mathrm{Fl}(v)\) is surjective.
\end{enumerate}

Let us then prove that the map \(\mathrm{Fl}(w)\) is surjective. Its image contains that of \(\mathrm{Fl}(\alpha')\) , because one has \(\alpha'\circ v=w\circ\alpha\) and the map \(\mathrm{Fl}(v)\) is surjective. Let b be a vertex of \(H'\) ; let a be a vertex of H such that b and \(u(a)\) are in the same orbit of \(H'\) and let f be an arrow of \(H'\) connecting \(u(a)\) to b. Then,

\[
h ^ {\prime} (u (a)) \cdot (\alpha^ {\prime} \circ \psi^ {\prime}) (f) = (\alpha^ {\prime} \circ \varphi^ {\prime}) (f) \cdot h ^ {\prime} (b),
\]

since \(h'\) is a homotopy connecting \(\alpha' \circ \varphi'\) to \(\alpha' \circ \psi'\) . The arrow \(h'(u(a)) = w(h(a))\) belongs to the image of \(\mathrm{Fl}(w)\) , as do the two arrows \(\alpha'(\psi'(f))\) and \(\alpha'(\varphi'(f))\) by what precedes. It follows that the arrow \(h'(b)\) belongs to the image of \(\mathrm{Fl}(w)\) , which proves that the image of \(\mathrm{Fl}(w)\) contains that of \(h'\) . By prop. 2 of II, p. 185, the map \(\mathrm{Fl}(w)\) is surjective.

Let \(g' \in \mathrm{Coh}(\varphi', \psi')_{u(a)}\) . Let g be an arrow of \(\mathrm{Coh}(\varphi, \psi)\) such that \(w(g) = g'\) . Denote by x and y the source and target of g; we have \(u(x) = u(y) = u(a)\) . Let \(g_1\) (resp. \(g_2\) ) be an arrow of G connecting a to x (resp. a to y) whose image under v is \(e_{u(a)}\) . Then, \(\alpha(g_1)g\alpha(g_2)^{-1}\) is an element of \(\mathrm{Coh}(\varphi, \psi)_a\) whose image under \(w_a\) is \(g'\) . Consequently, for every vertex a of G, the homomorphism \(w_a\) is surjective.

\begin{enumerate}
\def\labelenumi{\alph{enumi})}
\setcounter{enumi}{2}
\tightlist
\item
  Let us prove that, for every vertex \(a\) of G, the homomorphism \(w_{a}\) is injective. Considering successively the diagrams
\end{enumerate}

\[
\begin{array}{c c c} \mathrm{H} \xrightarrow {\varphi} \mathrm{G} & & \mathrm{H} \xrightarrow {v \circ \varphi} \mathrm{G} ^ {\prime} \\ \mathrm{Id} _ {\mathrm{H}} \Biggl \downarrow \quad \psi & \Biggl \downarrow v & u \Biggl \downarrow \quad \Biggl \downarrow \mathrm{Id} _ {\mathrm{G} ^ {\prime}} \\ \mathrm{H} \xrightarrow {v \circ \varphi} \mathrm{G} ^ {\prime} & & \mathrm{H} ^ {\prime} \xrightarrow {\varphi^ {\prime}} \mathrm{G} ^ {\prime} \\ & v \circ \psi & \psi^ {\prime} \end{array}
\]

We are then reduced to treating the following two cases: 1) We have H' = H and \(u = \mathrm{Id}_{\mathrm{H}}\) ; 2) We have G' = G and \(v = \mathrm{Id}_{\mathrm{G}}\) .

\begin{enumerate}
\def\labelenumi{\arabic{enumi})}
\tightlist
\item
  Suppose that H' = H and \(u = \mathrm{Id}_{\mathrm{H}}\) .
\end{enumerate}

Consider a morphism of groupoids \(v^{\prime}: G^{\prime} \to G\) which is inverse to v up to homotopy and a homotopy \(k: \operatorname{Som}(G) \to \operatorname{Fl}(G)\) connecting \(v^{\prime} \circ v\) to \(\mathrm{Id}_{\mathrm{G}}\) . By Remark 1 of II, p. 181, the maps \(\alpha \circ k \circ \varphi\) and \(\alpha \circ k \circ \psi\) are homotopies connecting respectively \(\alpha \circ v^{\prime} \circ \varphi^{\prime}\) to \(\alpha \circ \varphi\) and \(\alpha\circ v^{\prime}\circ\psi^{\prime}\) to \(\alpha\circ\psi\) . Consequently (cf. II, p. 180), the map

\[
\begin{array}{c} h _ {1} \colon \operatorname{Som} (\mathrm{H}) \to \operatorname{Fl} (\operatorname{Coh} (\varphi , \psi)), \\ x \mapsto (\alpha \circ k \circ \varphi) (x) \cdot h (x) \cdot ((\alpha \circ k \circ \psi) (x)) ^ {- 1} \end{array}
\]

is a homotopy connecting \(\alpha \circ v' \circ \varphi'\) to \(\alpha \circ v' \circ \psi'\) . By the universal property of cohomotopers (II, p. 185, prop. 3), there exists a unique morphism of groupoids \(w'\colon \mathrm{Coh}(\varphi', \psi') \to \mathrm{Coh}(\varphi, \psi)\) such that \(\alpha \circ v' = w' \circ \alpha'\) and \(h_1 = \mathrm{Fl}(w') \circ h'\) .

\[
\begin{array}{c} \mathrm{H} \xrightarrow [ \psi ]{\varphi} \mathrm{G} \xrightarrow {\alpha} \mathrm{Coh} (\varphi , \psi) \\ \mathrm{Id} _ {\mathrm{H}} \Biggl \downarrow \qquad \qquad \qquad \Biggl \downarrow v \qquad \qquad \Biggl \downarrow w \\ \mathrm{H} \xrightarrow [ \psi^ {\prime} ]{\varphi^ {\prime}} \mathrm{G} ^ {\prime} \xrightarrow {\alpha^ {\prime}} \mathrm{Coh} (\varphi^ {\prime}, \psi^ {\prime}) \\ \mathrm{Id} _ {\mathrm{H}} \Biggl \downarrow \qquad \qquad \qquad \Biggl \downarrow v ^ {\prime} \qquad \qquad \Biggl \downarrow w ^ {\prime} \\ \mathrm{H} \xrightarrow [ \psi ]{\varphi} \mathrm{G} \xrightarrow {\alpha} \mathrm{Coh} (\varphi , \psi)  . \end{array}
\]

In particular, we have

\[
\alpha \circ v ^ {\prime} \circ v = w ^ {\prime} \circ \alpha^ {\prime} \circ v = w ^ {\prime} \circ w \circ \alpha .
\]

Since k is a homotopy connecting \(v'\circ v\) to \(\mathrm{Id}_{\mathrm{G}}\) , \(\alpha\circ k\) is a homotopy connecting \(w'\circ w\circ\alpha\) to \(\alpha\) . Since \(\mathrm{Fl}(w'\circ w)\circ h=\mathrm{Fl}(w')\circ h'=h_{1}\) , we have, for every vertex x of H,

\[
\operatorname{Fl} \left(w ^ {\prime} \circ w\right) \circ h (x) \cdot (\alpha \circ k \circ \psi) (x) = h _ {1} (x) \cdot (\alpha \circ k \circ \psi) (x) = (\alpha \circ k \circ \varphi) (x) \cdot h (x),
\]

by definition of \(h_1\) . By prop. 5 of II, p. 186, applied to the morphisms of groupoids \(w' \circ w\) and \(\mathrm{Id}_{\mathrm{Coh}(\varphi, \psi)}\) , the map \(\alpha \circ k\) is a homotopy connecting \(w' \circ w\) to \(\mathrm{Id}_{\mathrm{G}}\) . In particular, \(w' \circ w\) is a homotopism.

For every vertex \(a\) of G, the group homomorphism \((w' \circ w)_a\) is therefore bijective (II, p. 182, prop. 1). It follows that the homomorphism \(w_a\) is injective, whence the result in case A).

\begin{enumerate}
\def\labelenumi{\arabic{enumi})}
\setcounter{enumi}{1}
\tightlist
\item
  Suppose we have G' = G and \(v = \mathrm{Id}_{\mathrm{G}}\) .
\end{enumerate}

Let \(x\) be a vertex of H'. The map \(\operatorname{Orb}(u)\) being surjective, there exists a vertex \(a\) of H and an arrow \(f\) of H' connecting \(u(a)\) to \(x\) . The arrows \(\alpha(\varphi'(f))^{-1}\) , \(h(a)\) and \(\alpha(\psi'(f))\) connect respectively \(\varphi'(x)\) to \(\varphi'(u(a)) = \varphi(a)\) , \(\varphi(a)\) to \(\psi(a)\) and \(\psi'(u(a)) = \psi(a)\) to \(\psi'(x)\) ; they are therefore composable in \(\operatorname{Coh}(\varphi, \psi)\) . Set

\[
h _ {2} (x) = \alpha (\varphi^ {\prime} (f)) ^ {- 1} \cdot h (a) \cdot \alpha (\psi^ {\prime} (f)).
\]

Let us verify that the arrow \(h_{2}(x)\) thus defined does not depend on the chosen elements a and f. Let \(a'\) be a vertex of H and let \(f'\) be an arrow of \(H'\) connecting \(u(a')\) to x. Since the map \(\operatorname{Orb}(u)\) is injective, the vertices a and \(a'\) of H belong to the same orbit and there exists an arrow \(c \in \operatorname{Fl}(\operatorname{H})\) connecting a to \(a'\) . Then, \(u(c)f'f^{-1}\) is a loop at \(u(a)\) in \(H'\) . By hypothesis (iii), there exists a vertex b of H such that the homomorphism \(u_{b}\) is surjective and an arrow \(c'\) of H from b to a; then, \(\operatorname{Int}(u(c'))(u(c)f'f^{-1})\) is a loop at \(u(b)\) in H, and it is therefore the image under \(u_{b}\) of a loop \(c''\) at b. Consequently, the arrow \(g = (c'c)^{-1}c''c'\) of H connects the vertex \(a'\) to the vertex a and satisfies \(f' = u(g)f\) . One then has

\[
\begin{array}{r l} & {\alpha (\varphi^ {\prime} (f ^ {\prime})) ^ {- 1} h (a ^ {\prime}) \alpha (\psi^ {\prime} (f ^ {\prime}))} \\ & {\qquad = \alpha (\varphi^ {\prime} (f)) ^ {- 1} \alpha (\varphi^ {\prime} (u (g))) ^ {- 1} h (a ^ {\prime}) \alpha (\psi^ {\prime} (u (g))) \alpha (\psi^ {\prime} (f))} \\ & {\qquad = \alpha (\varphi^ {\prime} (f)) ^ {- 1} \alpha (\varphi (g)) ^ {- 1} h (a ^ {\prime}) \alpha (\psi (g)) \alpha (\psi^ {\prime} (f))} \\ & {\qquad = \alpha (\varphi^ {\prime} (f)) ^ {- 1} \cdot h (a) \cdot \alpha (\psi^ {\prime} (f))} \end{array}
\]

since h is a homotopy connecting \(\alpha \circ \varphi\) to \(\alpha \circ \psi\) . This proves the announced independence.

By construction, we have \(h_{2}(u(x)) = h(x)\) for all \(x \in \operatorname{Som}(\mathrm{H})\) . We have thus defined a map \(h_{2}\) from \(\operatorname{Som}(\mathrm{H}')\) to \(\operatorname{Fl}(\operatorname{Coh}(\varphi, \psi))\) . Let c be an arrow of \(H'\) , denote its source by x and its target by y. Let a be a vertex of H and let f be an arrow of \(H'\) connecting \(u(a)\) to x. Then fc is an arrow of \(H'\) which connects \(u(a)\) to y. By definition of \(h_{2}\) , we thus have

\[
\begin{array}{r l} & h _ {2} (x) \alpha (\psi^ {\prime} (c)) = \alpha (\varphi^ {\prime} (f)) ^ {- 1} h (a) \alpha (\psi^ {\prime} (f)) \alpha (\psi^ {\prime} (c)) \\ & \qquad = \alpha (\varphi^ {\prime} (c)) \alpha (\varphi^ {\prime} (f c)) ^ {- 1} h (a) \alpha (\psi^ {\prime} (f c)) \\ & \qquad = \alpha (\varphi^ {\prime} (c)) h _ {2} (y). \end{array}
\]

This proves that \(h_{2}\) is a homotopy connecting \(\alpha \circ \varphi'\) to \(\alpha \circ \psi'\) .

By the universal property of cohomotopers, there exists a unique groupoid morphism \(w'\colon \mathrm{Coh}(\varphi', \psi') \to \mathrm{Coh}(\varphi, \psi)\) such that \(w' \circ \alpha' = \alpha\) and \(h_2 = \mathrm{Fl}(w') \circ h'\) . We have \(w' \circ w \circ \alpha = w' \circ \alpha' = \alpha\) and \(\mathrm{Fl}(w' \circ w) \circ h = \mathrm{Fl}(w') \circ h' \circ \mathrm{Som}(u) = h_2 \circ u = h\) by definition of \(h_2\) . By the universal property of cohomotopers, this implies that \(w' \circ w = \mathrm{Id}_{\mathrm{Coh}(\varphi, \psi)}\) . In particular, for every \(a \in \mathrm{Som}(G)\) , the homomorphism \(w_a\) is injective.

It then follows from prop. 1 of II, p. 182 that the morphism w is a homotopism, whence the theorem.
% semantic-fragments: legacy-input-seam
\relax{}
